\begin{frame}{Актуальность моделирования ECAP}
\begin{itemize}
\item Вызванный составной потенциал действия (ECAP) характеризует электрический ответ возбуждённых нервных волокон при стимуляции спинного мозга.
\item Амплитуда ответа зависит от интенсивности и длительности импульса, геометрии электродов и условий регистрации.
\item Артефакт стимуляции и шум затрудняют выделение нейронного сигнала.
\item Биофизическая модель позволяет исследовать вклад отдельных механизмов и условия сопоставимости измерений.
\end{itemize}
\SlideGap
\textbf{Основная целевая величина:} кривая роста ECAP --- зависимость амплитуды ответа от интенсивности стимуляции.
\SourceNote{Anaya et al., Neuromodulation, 2020. DOI: \href{https://doi.org/10.1111/ner.12965}{10.1111/ner.12965}. Brucker-Hahn et al., Journal of Neural Engineering, 2023. DOI: \href{https://doi.org/10.1088/1741-2552/aceca4}{10.1088/1741-2552/aceca4}.}
\end{frame}
